\section{引言}
格点 QCD 中从欧几里得空间关联函数的蒙特卡罗估计提取强子质量与矩阵元时，
如果所用算符对目标态耦合更强、对较高能级的污染态耦合更弱，则提取会更加
可靠和精确。对于包含胶子的态，在格点 QCD 中构造此类算符的一个关键要素是
链接变量涂抹（smearing）。由涂抹过或 fuzz 过的链接构造的算符，与理论高频
模式的混合被大幅压低。已经证明，此类算符的使用对胶球谱\cite{colin2}、
杂化介子质量\cite{hybrids,jkm1}、torelon 谱\cite{torelons}以及静态
夸克-反夸克势的激发谱\cite{jkm2}的确定尤为有益。

链接变量光滑化在改进格点作用量的构造中也正发挥着日益重要的作用。
文献~\cite{fat1} 表明，在 staggered 夸克作用量中使用所谓的胖链接
（fat links）能显著减小味对称性破缺。随后，人们用涂抹过的链接\cite{fat2}
构造了具有更好转动不变性的超立方费米子作用量。采用称为超立方（HYP）胖链接
的链接涂抹变换\cite{hyp} 的 staggered 费米子作用量，其味对称性相对标准
作用量提高了一个数量级。所谓 Asqtad 改进 staggered 夸克作用量%
\cite{asqtad1,asqtad2,asqtad3} 也利用链接加粗来减小味对称性破缺。
另一类利用涂抹链接变量的费米子作用量是胖链接无关 clover（FLIC）作用量%
\cite{flic}。FLIC 费米子属 Wilson 型，其作用量包含一个用涂抹链接构造的
无关 clover 改进项。胖链接还被用于借助近似重整化群变换构造离散误差更小的
规范作用量\cite{fatgauge}。

在胶子算符构造中最常用的链接 fuzz 算法见文献~\cite{APEsmear}：将格点上
每一条空间链接 $U_j(x)$ 替换为其自身加上实权重 $\rho$ 乘以其四个相邻
（空间）staple 组合之和，再投影回 $SU(3)$。如此将该 fuzz 步骤迭代
$n_\rho$ 次即得最终的 fuzz 链接变量。经验上发现，投影回 $SU(3)$ 是涂抹
的关键要素；不作此投影的链接涂抹方法效果差得多。到 $SU(3)$ 的投影并不
唯一，必须仔细定义以保持链接变量的全部对称性质。文献中出现了多种实现该
投影的办法。给定 $3\times 3$ 矩阵 $V$，其到 $SU(3)$ 的投影 $U$ 常取为使
$\ReTr(UV^\dagger)$ 最大的矩阵 $U\in SU(3)$；完成这种最大化需要迭代
程序。另一种做法\cite{su3projectA}是把投影后的矩阵 $U$ 定义为
$ U = V\ (V^\dagger V)^{-1/2} \det(V^{-1}V^\dagger)^{1/6}.$

尽管上述涂抹流程在实践中行之有效，这一投影却多少令人不快，因为它可被视
为停留在群内的一种人为而突兀的方式。更重要的是，投影固有的分支切割与不可
微性会妨碍、甚至使现代蒙特卡罗更新技术无法应用，例如混合蒙特卡罗
（HMC）\cite{hmc}，这类方法需要知道作用量对某条链接变量微小变化的响应。

本文提出并检验了一种规避上述问题的链接涂抹方法。该方法在整个有限复平面
上处处解析，并以保持在 $SU(3)$ 内的方式利用指数函数，从而免去任何投影回
群的步骤。由于这一构造，该算法适用于任意李群。第~\ref{sec:analytic} 节
描述该方法；在构造算符时以及在计算作用量对链接变量变化的响应时的实际实现
分别于第~\ref{sec:implement} 节与第~\ref{sec:force} 节详细给出；
数值检验在第~\ref{sec:tests} 节给出。借助涂抹方格值以及与静态夸克-反夸克
势相联系的有效质量 $aE_{\rm eff}(t)$，我们观察到与现行标准链接涂抹方法
相比效果并无退化，只是发现对涂抹参数的敏感性增大。结果还表明，由涂抹链接
变量构造的格点作用量与算符受辐射修正的影响可能要小得多，因为通常占主导
地位的大蝌蚪图贡献被急剧削减。
